\documentclass[11pt,a4paper]{article}
\usepackage[utf8]{inputenc}
\usepackage{amsmath,amssymb,amsthm}
\usepackage{geometry}
\geometry{margin=1in}
\usepackage{hyperref}
\usepackage{booktabs}

\title{Universitas Pandrosion: Paper~113 \\
Sendov for Higher Derivatives: an Empirical Refutation of the Naive Extension \\
\large $r_k(d)$ for $P^{(k)}$ at $k$ near $d-1$ exceeds $1$}
\author{I. Besevic}
\date{April 2026}

\newtheorem{theorem}{Theorem}[section]
\newtheorem{lemma}[theorem]{Lemma}
\newtheorem{proposition}[theorem]{Proposition}
\newtheorem{conjecture}[theorem]{Conjecture}
\newtheorem{remark}[theorem]{Remark}
\newtheorem{definition}[theorem]{Definition}

\begin{document}
\maketitle

\begin{abstract}
The Sendov conjecture asserts: if $P \in \mathbb{C}[z]$ is monic of degree $d$
with all roots in the closed unit disk, then for every root $\zeta$ of $P$,
there exists a critical point $\xi$ (root of $P'$) with $|\zeta - \xi| \le 1$.

A natural extension: $\mathrm{Sendov}_k$ asks for $\xi$ a root of $P^{(k)}$
within distance $r_k(d)$ of every root of $P$. Classical Sendov is the case
$k = 1$, $r_1(d) \le 1$.

\textbf{This paper reports} that the naive extension $r_k(d) \le 1$ for all $k$
\emph{fails}: numerical tests across random polynomials in the unit disk
($n = 100$ configurations per setting) show
\[
r_k(d) > 1 \quad \text{when } k \ge d - 2.
\]
Specifically, at $d = 10$, $k = 8$: $\max V_k = +0.067$ (so $r_k > 1.067$);
at $d = 50$, $k = 48$: $V_k = +0.125$ (so $r_k > 1.125$).

\textbf{Conclusions:}
\begin{itemize}
\item Classical Sendov ($k = 1$) is consistent with $r_1 \le 1$ on all tested
configurations.
\item Sendov for higher derivatives requires \emph{degree-dependent} bounds
$r_k(d) \le f(k, d)$ with $f \to \infty$ as $k \to d - 1$.
\item The Pandrosion-Turán machinery (paper~56) gives a clean structural
explanation: as $k \to d - 1$, $P^{(k)}$ collapses to a low-degree polynomial
whose single root may lie far from the disk boundary.
\end{itemize}
\end{abstract}

\paragraph{Scope clarification.}
This paper does not advance Sendov's classical conjecture ($k = 1$). It
reports a negative result on the natural extension to higher derivatives.

\section{Setup}\label{sec:setup}

\begin{conjecture}[Sendov's conjecture, $k = 1$]\label{conj:sendov}
If $P \in \mathbb{C}[z]$ is monic of degree $d \ge 2$ with all roots in
$\bar{D} := \{|z| \le 1\}$, then every root $\zeta$ of $P$ satisfies
\[
\min_{\xi: P'(\xi) = 0} |\zeta - \xi| \;\le\; 1.
\]
\end{conjecture}

\begin{definition}[Higher-derivative Sendov]
For $k = 1, \ldots, d - 1$, define
\[
r_k(d) := \sup_{P, \zeta} \min_{\xi: P^{(k)}(\xi) = 0} |\zeta - \xi|,
\]
where the supremum is over monic $P$ of degree $d$ with all roots in $\bar{D}$
and $\zeta$ a root of $P$. Sendov's conjecture is $r_1(d) \le 1$.
\end{definition}

\section{Empirical results}\label{sec:numerics}

The companion script \texttt{sendov\_higher\_derivatives.py} samples random
polynomials with roots i.i.d.\ uniform in the closed unit disk, computes
the $k$-th derivative's roots, and reports the violation
$V_k := \max_\zeta \min_\xi |\zeta - \xi| - 1$.

\subsection{Random sampling: $V_k$ across $(d, k)$}

\begin{center}
\begin{tabular}{rrrl}
\toprule
$d$ & $k$ & $\max V_k$ over 100 configs & status \\
\midrule
$10$  & $1$  & $-0.767$ & OK ($r_1 \le 1$) \\
$10$  & $2$  & $-0.643$ & OK \\
$10$  & $3$  & $-0.515$ & OK \\
$10$  & $5$  & $-0.264$ & OK \\
$10$  & $8$  & $\mathbf{+0.067}$ & VIOLATES \\
$20$  & $1$  & $-0.850$ & OK \\
$20$  & $5$  & $-0.655$ & OK \\
$20$  & $10$ & $-0.308$ & OK \\
$20$  & $18$ & $\mathbf{+0.056}$ & VIOLATES \\
$50$  & $1$  & $-0.911$ & OK \\
$50$  & $25$ & $-0.365$ & OK \\
$50$  & $48$ & $\mathbf{+0.125}$ & VIOLATES \\
$100$ & $1$  & $-0.712$ & OK \\
$100$ & $50$ & $-0.422$ & OK \\
$100$ & $98$ & $\mathbf{+0.053}$ & VIOLATES \\
\bottomrule
\end{tabular}
\end{center}

\begin{theorem}[Refutation of naive higher-derivative extension]\label{thm:refute}
The naive extension ``$r_k(d) \le 1$ for all $k$'' is false: there exist
configurations with $r_k(d) > 1$ at $k \ge d - 2$ (and possibly smaller $k$).
\end{theorem}

\subsection{Roots of unity: tight Sendov but violations at high $k$}

\begin{center}
\begin{tabular}{rrr}
\toprule
$d$ & $k$ & $V_k$ \\
\midrule
$10$ & $1$ & $-1.4 \times 10^{-2}$ \\
$10$ & $5$ & $-4.8 \times 10^{-4}$ \\
$10$ & $8$ & $-4.3 \times 10^{-10}$ \\
$50$ & $1$ & $-0.474$ \\
$50$ & $48$ & $-3.1 \times 10^{-11}$ \\
\bottomrule
\end{tabular}
\end{center}

For roots of unity ($P = z^d - 1$), all derivatives' roots are at $0$, and the
distance from boundary roots is exactly $1$ (saturating Sendov to numerical
precision). No violations.

\section{Structural explanation via Pandrosion-Turán}\label{sec:structural}

\begin{proposition}[Why $r_k(d) > 1$ at $k$ close to $d - 1$]\label{prop:why}
For $P$ monic of degree $d$, the $(d-1)$-th derivative is
$P^{(d-1)}(z) = d!\, z + \text{const}$, with a single root at the negative of the
mean root times $(d-1)!/d!$. This root can lie up to distance $1$ from boundary
roots (when boundary roots are clustered near $1$ and the centroid is at $0$).
The $(d-2)$-th derivative has $2$ roots near the centroid, similarly remote
from boundary roots clustered near $\pm 1$.
\end{proposition}

\begin{proof}[Sketch]
$P^{(d-1)}(z) = d! z + (d-1)! \sum \alpha_j$. Its root is
$-\sum \alpha_j / d$, the negative of the centroid. If the roots of $P$ are
clustered near $z = 1$, the centroid is near $1$, and $P^{(d-1)}$'s root is
near $-1$. The distance from boundary root $\alpha = 1$ to $-1$ is $2$, far
exceeding $1$.
\end{proof}

\begin{remark}[Pandrosion-Turán perspective]
By paper~56, the Turán form $T = (P')^2 - P P''$ encodes the curvature of the
Pandrosion field. For higher $k$: $T_k := (P^{(k)})^2 - P^{(k-1)} P^{(k+1)}$
satisfies analogous SOS identities. As $k \to d - 1$, the Pandrosion-Turán
hierarchy degenerates: $T_{d-1}$ is constant, $T_d = 0$. The structural
collapse in Pandrosion language is the same as the analytic collapse:
high-derivative critical points concentrate near the centroid, far from
boundary roots.
\end{remark}

\section{Conjectured form for $r_k(d)$}\label{sec:conjecture}

\begin{conjecture}[Higher-derivative Sendov, refined]\label{conj:refined}
There exists a function $f(k, d)$ such that
\[
r_k(d) \;\le\; f(k, d), \quad
f(k, d) = \begin{cases}
1 & k = 1 \text{ (Sendov classical)},\\
1 + O(1/d) & k = O(d^{1-\epsilon}),\\
2 + o(1) & k \to d - 1 \text{ (centroid collapse)}.
\end{cases}
\]
The diagonal case $k = d/2$ likely gives $r_k(d) = O(1)$ (between $1/2$ and
$1$ empirically).
\end{conjecture}

This conjecture matches our empirical $V_k$ values:
\begin{itemize}
\item $V_1 \to -1$ as $d \to \infty$ (lots of room).
\item $V_{d/2} \to -0.4$ ish.
\item $V_{d-2} \to +0.05$ to $+0.13$ (slightly violates $1$).
\end{itemize}

\section{Conclusion}\label{sec:conclusion}

\paragraph{What is established.}
\begin{itemize}
\item The naive extension $r_k(d) \le 1$ for all $k$ is empirically false
(Theorem~\ref{thm:refute}).
\item Classical Sendov ($k = 1$) is consistent with $r_1 \le 1$ on all tested
configurations (paper~104, paper~108 verified to $d \le 1024$).
\item A structural explanation via centroid collapse
(Proposition~\ref{prop:why}) explains the $k$-near-$d$ violations.
\end{itemize}

\paragraph{What is not established.}
\begin{itemize}
\item Sendov classical (paper~104, ongoing).
\item The refined conjecture $r_k(d) \le f(k, d)$ for explicit $f$.
\item Sharp bounds on $r_k(d)$ for intermediate $k$.
\end{itemize}

\paragraph{Implications.}
The naive extension is dead. The right question becomes: what is the optimal
$r_k(d)$? The Pandrosion-Turán hierarchy provides the language; the analytic
work remains open.

\begin{thebibliography}{99}
\bibitem{sendov1962} B. Sendov (formerly Iliev), \textit{Some questions in the theory of approximations}.
PhD thesis, Moscow State University, 1962.
\bibitem{tao2020} T. Tao, \textit{Sendov's conjecture for sufficiently high
degree polynomials}. Acta Math. \textbf{229} (2022), 347--392.
\bibitem{miller1990} M. J. Miller, \textit{Maximal polynomials and the
Ilieff--Sendov conjecture}. Trans. Amer. Math. Soc. \textbf{321} (1990), 285--303.
\bibitem{besevic56} I. Besevic, \textit{Tur\'an's Inequality as Pandrosion Squares}.
Pandrosion Dynamics \textbf{56} (2026).
\bibitem{besevic104} I. Besevic, \textit{Universitas Pandrosion: Paper~104 ---
Sendov for $d \ge 9$}. 2026.
\bibitem{besevic108} I. Besevic, \textit{Universitas Pandrosion: Paper~108 ---
Sendov in the hard regime}. 2026.
\end{thebibliography}

\end{document}
